\reading{IM-10}{The Rise and Risks of Private Credit}{DEFER}{0}

\srcnote{``The Last Mile: Financial Vulnerabilities and Risks, Chapter 2: The
Rise and Risks of Private Credit'', IMF Global Financial Stability Report (April
2024). No GARP source book covers this reading; gap-fills are one rung below a
GARP source.}

\begin{imkeybox}[title={The four objectives in this reading}]
\begin{enumerate}[label=\textbf{\alph*.}]
\item Describe characteristics of private credit, including its typical investors and borrowers, and compare private credit to other types of loans and fixed-income instruments.
\item Explain the return profile and growth profile of the private credit asset class, and compare the historical returns of private credit to those of other asset classes.
\item Describe and assess the risks and vulnerabilities related to private credit, and explain how private credit can pose risks to financial stability.
\item Assess potential policy recommendations that could help mitigate the risks associated with private credit.
\end{enumerate}
\end{imkeybox}

\begin{imtrapbox}[breakable=false,title={The second unlabelled objective}]
The source notes label a, b and c and stop. IM-10 d is nowhere in the notes as a
heading --- but its content is present throughout, split across six ``How to
mitigate it'' blocks attached to the six risk categories. Like IM-9 c, this reads
as a gap and is not one. It is gathered under IM-10 d below.
\end{imtrapbox}

% ------------------------------------------------------------------
\lo{IM-10 a}{Describe characteristics of private credit, including its typical investors and borrowers, and compare private credit to other types of loans and fixed-income instruments.}

\begin{imdefbox}[title={Private credit}]
\term{Nonbank corporate credit} provided through direct, bilateral agreements
between lenders and borrowers. It excludes public securities and traditional
commercial bank lending, and primarily serves \term{middle-market companies} ---
those with annual revenues between \$100 million and \$1 billion.
\end{imdefbox}

\begin{itemize}
\item \term{Direct lending structure.} Loans established directly between asset
managers and borrowers, without commercial banks or public markets.
\item \term{Customized loan terms.} Deferred interest payments, bespoke covenants,
flexible repayment schedules.
\item \term{Higher costs.} Interest payments and total borrowing costs are
typically higher than traditional bank loans, reflecting higher credit risk and an
\term{illiquidity premium}.
\item \term{Enhanced covenants.} Stricter limits on dividend payments, share
repurchases and asset sales, to protect lenders.
\end{itemize}

\renewcommand{\arraystretch}{1.2}
\noindent\begin{tabularx}{\textwidth}{@{}>{\raggedright\arraybackslash}p{2.4cm}X@{}}
\toprule
\textbf{\sffamily\small Investors} & Long-term institutional capital: \term{pension funds, insurance companies, sovereign wealth funds and family offices}, seeking higher yields and diversification beyond public markets. Structures used: \term{closed-end funds} (common, similar to private equity); \term{business development companies (BDCs)}, publicly traded vehicles offering retail access; and \term{collateralized loan obligations (CLOs)}, pooled loans securitized into tranches. \\
\textbf{\sffamily\small Borrowers} & Small to mid-sized companies, often in healthcare and information technology. They may operate with high leverage or weaker earnings and \emph{fail to meet the criteria for traditional bank financing}, yet still attract credit due to \term{private equity backing} and strong lender relationships. Loans are commonly \term{floating-rate} and require more collateral and covenants than bank loans. \\
\textbf{\sffamily\small Alternatives} & Traditional bank loans, public debt issuance, and \term{broadly syndicated loans}, which involve multiple lenders sharing risk for large or risky borrowers. A leveraged loan is a type of syndicated loan. \\
\bottomrule
\end{tabularx}
\renewcommand{\arraystretch}{1}
\figcap{Who lends, who borrows, and what else they could have done. The borrower
row carries the tension in the asset class: these are firms a bank would decline,
lent to at floating rates by investors whose own liabilities are long-dated and
fixed.}

% ------------------------------------------------------------------
\lo{IM-10 b}{Explain the return profile and growth profile of the private credit asset class, and compare the historical returns of private credit to those of other asset classes.}

Private credit has expanded rapidly since the 2007--2009 financial crisis, driven
by \term{low interest rates}, which made traditional financing less attractive for
investors, and \term{stricter regulations on banks}, which limited their ability
to lend to mid-sized, riskier companies. As banks reduced exposure, private credit
filled the gap, with insurance companies and other institutional investors
entering because of \term{lighter capital and risk requirements}.

\begin{imkeybox}[title={The size, as of 2023}]
Approximately \term{\$2.1 trillion} in loans and undeployed capital --- nearing
the size of the global \term{high-yield bond market}. Most activity is in North
America, with growth increasing in Europe and Asia.
\end{imkeybox}

Originally focused on mid-sized companies, private credit has expanded to serve
larger corporations that traditionally relied on syndicated loans or public bonds,
because of customized loan structures, \term{reduced disclosure requirements}, and
lower costs than public bond issuance.

\begin{imkeybox}[title={Historical returns}]
Private credit has had \term{lower losses than high-yield (junk) bonds} and
\term{similar losses to leveraged loans}. This is partly because many loans are
backed by private equity firms who support the borrower, and are secured by
collateral.
\end{imkeybox}

\begin{imtrapbox}[breakable=false,title={Is it cyclical? The evidence cuts both ways}]
Some evidence shows private credit lenders still provided funds in tough times ---
the March 2020 COVID crash, for instance. But funding availability is
\emph{somewhat} tied to the stock market: when equities do well, private credit
flows more freely, and when banks tighten their lending standards, private credit
lending tends to slow too. The curriculum does not settle this. An option that
states flatly that private credit is countercyclical, or flatly that it is
procyclical, has taken a side the source does not take.
\end{imtrapbox}

% ------------------------------------------------------------------
\lo{IM-10 c}{Describe and assess the risks and vulnerabilities related to private credit, and explain how private credit can pose risks to financial stability.}

There has been a shift from regulated lending --- banks and public bonds --- to
less regulated private credit, which involves \term{nonrated, nontraded loans with
minimal transparency}. That makes it hard to detect risks early, and there is
currently limited data to assess the full impact on financial stability. Risks
appear moderate, due to duration matching and low leverage, but concerns remain:

\begin{itemize}
\item Lack of transparency and regulation.
\item \term{No major stress test yet} --- there has been no significant market
correction.
\item Risk of \term{lower underwriting standards} due to competition between
funds.
\item Potential for reduced risk premiums and weaker covenants with rising supply.
\end{itemize}

\renewcommand{\arraystretch}{1.25}
\noindent\begin{tabularx}{\textwidth}{@{}>{\raggedright\arraybackslash}p{2.3cm}X@{}}
\toprule
\textbf{\sffamily\small Risk} & \textbf{\sffamily\small How it arises} \\
\midrule
Credit & Borrowers usually take \term{floating-rate} loans. If rates rise sharply their payments increase, making repayment harder; borrowers are often given temporary relief by \term{adding unpaid interest to the outstanding debt} to be repaid later. \\
Liquidity & Funds hold illiquid private corporate loans and manage this with \term{long-term lockups, limited redemption windows, gates and suspensions}. \term{Semiliquid evergreen structures} --- accepting ongoing capital and investing continuously --- carry higher liquidity risk, especially as they attract individual investors seeking partial liquidity. A further risk arises if borrowers \term{draw down credit lines simultaneously}, forcing capital calls that transfer the liquidity burden onto investors. \\
Leverage & Funds may appear to use limited leverage, but leverage exists \term{at the investor, fund and borrower levels at once}, and this layering is hard to monitor. Funds lever through SPVs and collateralized fund obligations; borrowers are often funded by both the credit fund and the backing PE firm. \\
Valuation & Valuations lack reliability: no comparable transactions, infrequent updates, no secondary market. Price changes are slow, so \term{stale valuations persist for quarters}. This permits overvaluation, lets informed investors exit early, \term{overstates returns} and \term{inflates manager fees}. \\
Inter-\\ connected & Many private credit managers are \emph{also} PE managers and borrowers are often PE-backed, which reduces credit risk but creates \term{conflict-of-interest} concerns. Banks' exposure is minimal ($<$1\% of assets) and secured; pension funds and insurers allocate \term{$\approx$3\%}. PE-influenced life insurers hold more illiquid assets yet have \term{weaker solvency ratios} than peers. \\
Underwriting & Competition with banks and bond markets for large borrowers has led to \term{weaker underwriting standards, underpriced loans and looser covenants}. In a downturn this could mean heavy losses for lenders. \\
\bottomrule
\end{tabularx}
\renewcommand{\arraystretch}{1}
\figcap{Six vulnerabilities, one common root. Every row traces back to the same
property --- these loans are not traded, so they are not priced, not rated and not
observed. Note that the leverage row is the only one where the risk is
\emph{layered}: the same underlying exposure is levered three times over by three
different parties, none of whom can see the other two.}

% ------------------------------------------------------------------
\lo{IM-10 d}{Assess potential policy recommendations that could help mitigate the risks associated with private credit.}

\srcnote{The objective the source notes leave unlabelled. The material below is
theirs, gathered from the six ``How to mitigate it'' blocks; only the heading is
supplied.}

\renewcommand{\arraystretch}{1.25}
\noindent\begin{tabularx}{\textwidth}{@{}>{\raggedright\arraybackslash}p{2.3cm}X@{}}
\toprule
\textbf{\sffamily\small Against} & \textbf{\sffamily\small Recommendation} \\
\midrule
Credit risk & Insurers and pension funds usually do not monitor loan credit quality directly, relying instead on \term{internal valuation models} and \term{credit rating agencies}. \\
Liquidity risk & Even where funds have tried to match liquidity better, \term{open-ended funds} remain a problem: if many unsophisticated investors panic and exit during stress, the fund may not be able to meet redemptions. \\
Leverage risk & Reporting is incomplete --- \term{SPVs, feeder funds and offshore structures hide true risk}, and inconsistencies across countries and sectors make oversight difficult. Banks lending to private credit firms should require \term{liquidity stress testing}, preparation for \term{asset writedowns}, and testing of exposure to \term{borrower credit lines}. Data deficiencies call for better cross-border reporting rules, regulatory coordination to spot hidden leverage, and \term{leverage caps} if risk becomes systemically dangerous. \\
Inter-\\ connected & \term{Regulatory arbitrage} --- entities seeking favourable regulation locally or globally --- may increase risk, particularly among insurers and pension funds with substantial holdings. Reducing it requires \term{harmonising risk assessment across sectors and borders}, and international cooperation between regulators. \\
Conduct risk & As more unsophisticated investors enter, conflicts of interest arise, especially where funds control \emph{both sides} of a transaction. Supervisors must increase transparency, prioritise investor interests, ensure fairness in sales practices and suitability, and put controls in place so retail investors understand credit and liquidity risks and \term{redemption restrictions}. \\
\bottomrule
\end{tabularx}
\renewcommand{\arraystretch}{1}
\figcap{The policy response, read off the risks. Note what is \emph{not} proposed:
nowhere does the curriculum recommend prohibiting the activity or forcing private
credit into public markets. Every recommendation is about \term{seeing} the risk
--- reporting, coordination, transparency, disclosure --- with leverage caps named
only as a contingency if risk becomes systemically dangerous.}

\begin{imtrapbox}[breakable=false,title={The conflict that matters most}]
Frequent redemptions may incentivise fund managers to \term{manipulate valuations}
or impose conditions on redemptions to retain investor funds. That links the
valuation risk and the conduct risk into one mechanism: stale valuations are not
merely a measurement problem, they are something a manager has a reason to
prefer.
\end{imtrapbox}

\begin{imtrapbox}[breakable=false,title={Consolidated trap summary --- IM-10}]
\begin{itemize}
\item \term{Definition, IM-10 a.} Middle market is \$100 million to \$1 billion in
\emph{revenue}, not in assets or loan size.
\item \term{Number, IM-10 b.} \$2.1 trillion as of 2023, comparable to the global
\emph{high-yield bond} market --- not to the leveraged loan or investment-grade
market.
\item \term{Polarity, IM-10 b.} Losses have been \emph{lower} than high-yield
bonds and \emph{similar} to leveraged loans. Two different comparisons in one
sentence.
\item \term{Scope, IM-10 b.} The curriculum does not declare private credit
cyclical or countercyclical. Evidence points both ways.
\item \term{Sibling, IM-10 c.} Leverage risk is the \emph{layered} one --- investor,
fund and borrower. Interconnectedness is about overlapping \emph{parties}, not
overlapping borrowings.
\item \term{Number, IM-10 c.} Banks' exposure is under 1\% of assets; pension funds
and insurers allocate around 3\%.
\item \term{Role, IM-10 c.} PE-influenced life insurers hold \emph{more} illiquid
assets and have \emph{weaker} solvency ratios than peers. Both halves point the
same way.
\item \term{Scope, IM-10 d.} The recommendations are about visibility, not
prohibition. Leverage caps are a contingency, not a proposal.
\end{itemize}
\end{imtrapbox}

% ==================================================================
